% ============================================================================
% Feed-Forward Networks
% Included immediately after the corresponding \section command in the
% architecture chapter. This file intentionally contains no \chapter,
% \section, package imports, or document-level commands.
% ============================================================================

Positional encoding supplies token representations with information about
their positions in a sequence. Position information, however, does not by
itself specify how the features represented by each token are transformed
into increasingly useful internal representations. A Transformer therefore
requires computational components that operate on the feature dimensions of
token representations in addition to components that establish relationships
between different token positions. The feed-forward network (FFN) provides
this feature-transformation mechanism. In a conventional Transformer block,
the FFN is applied independently to each token representation, using the
same learned parameters at every sequence position
\parencite{dai2024deepseekmoe,cai2025survey}.

The distinction between position-wise processing and sequence-wise
interaction is fundamental. An FFN transforms the feature vector associated
with one token without directly combining it with the feature vectors of
other tokens. The parameters are shared across positions, so the same
nonlinear mapping is applied throughout the sequence. This arrangement
complements the self-attention sublayer, which supplies interactions between
token positions. The present section develops the mathematical structure of
the FFN, examines its parameter and computational requirements, and then
follows the architectural transition from dense FFNs to sparsely activated
mixtures of FFN experts.

\subsection{Position-Wise Feed-Forward Networks}
\label{subsec:ffn-position-wise}

Let a sequence contain \(T\) tokens, and let
\(d_{\mathrm{model}}\) denote the dimensionality of each token
representation. The representation of token \(t\) is a vector
\(\mathbf{x}_t\in\mathbb{R}^{d_{\mathrm{model}}}\). A sequence can be
represented by a matrix
\[
    \mathbf{X}
    =
    \begin{bmatrix}
        \mathbf{x}_1 & \mathbf{x}_2 & \cdots & \mathbf{x}_T
    \end{bmatrix}
    \in
    \mathbb{R}^{d_{\mathrm{model}}\times T},
\]
where each column corresponds to one token. This column-wise convention
will be used throughout the mathematical development.

A standard two-projection FFN first maps the input representation into an
intermediate feature space, applies an element-wise nonlinear activation,
and then maps the result back into the model hidden space. For one token,
the computation is
\begin{align}
    \mathbf{a}_t
    &=
    \mathbf{W}_1\mathbf{x}_t+\mathbf{b}_1,
    \label{eq:ffn-first-projection}\\
    \widetilde{\mathbf{a}}_t
    &=
    \phi(\mathbf{a}_t),
    \label{eq:ffn-elementwise-activation}\\
    \mathbf{y}_t
    &=
    \mathbf{W}_2\widetilde{\mathbf{a}}_t+\mathbf{b}_2.
    \label{eq:ffn-second-projection}
\end{align}
Here, \(d_{\mathrm{ff}}\) is the intermediate hidden dimension,
\(\mathbf{W}_1\in\mathbb{R}^{d_{\mathrm{ff}}\times d_{\mathrm{model}}}\)
is the input-projection matrix,
\(\mathbf{b}_1\in\mathbb{R}^{d_{\mathrm{ff}}}\) is its bias vector,
\(\mathbf{W}_2\in\mathbb{R}^{d_{\mathrm{model}}\times d_{\mathrm{ff}}}\)
is the output-projection matrix, and
\(\mathbf{b}_2\in\mathbb{R}^{d_{\mathrm{model}}}\) is the output bias.
The activation function \(\phi\) acts element-wise on its vector argument.
The dimensions therefore satisfy
\[
\begin{aligned}
    \mathbf{x}_t
        &\in \mathbb{R}^{d_{\mathrm{model}}},&
    \mathbf{a}_t,\widetilde{\mathbf{a}}_t
        &\in \mathbb{R}^{d_{\mathrm{ff}}},&
    \mathbf{y}_t
        &\in \mathbb{R}^{d_{\mathrm{model}}}.
\end{aligned}
\]
For a bias-free formulation, the bias vectors are omitted. The essential
structure is the same: an input projection, an element-wise nonlinearity,
and an output projection. The conventional FFN is consequently a
two-layer multilayer perceptron whose input and output dimensions agree,
while its intermediate dimension may be wider. The description of standard
FFN experts as linear--activation--linear networks is also used in the
MoE literature \parencite{dai2024deepseekmoe,cai2025survey}.

The complete sequence is processed by applying the same matrices and bias
vectors to every column:
\begin{align}
    \mathbf{A}
    &=
    \mathbf{W}_1\mathbf{X}
    +\mathbf{b}_1\mathbf{1}_T^{\top},
    \label{eq:ffn-batch-first}\\
    \widetilde{\mathbf{A}}
    &=
    \phi(\mathbf{A}),
    \label{eq:ffn-batch-activation}\\
    \mathbf{Y}
    &=
    \mathbf{W}_2\widetilde{\mathbf{A}}
    +\mathbf{b}_2\mathbf{1}_T^{\top},
    \label{eq:ffn-batch-second}
\end{align}
where \(\mathbf{1}_T\in\mathbb{R}^{T}\) is the all-ones vector. Thus,
\(\mathbf{A},\widetilde{\mathbf{A}}\in
\mathbb{R}^{d_{\mathrm{ff}}\times T}\) and
\(\mathbf{Y}\in\mathbb{R}^{d_{\mathrm{model}}\times T}\).
The matrix multiplication acts on the feature dimension of each column;
there is no multiplication between different token columns. Therefore,
for every \(t\), the output \(\mathbf{y}_t\) depends on
\(\mathbf{x}_t\) and the shared FFN parameters, but not directly on
\(\mathbf{x}_{s}\) for \(s\ne t\).

The nonlinearity is necessary for the composition to represent a
nonlinear mapping. If \(\phi\) were the identity function and the biases
were zero, then
\[
    \mathbf{y}_t
    =
    \mathbf{W}_2\mathbf{W}_1\mathbf{x}_t.
\]
The product \(\mathbf{W}_2\mathbf{W}_1\) is itself a matrix, so the two
projections would collapse to one linear transformation. An intervening
nonlinear activation prevents this reduction in general and allows the
FFN to implement a nonlinear feature transformation. The specific
mathematical properties of the activation function are treated in the
subsequent section on activation functions.

\subsection{Parameterization and Computational Complexity}
\label{subsec:ffn-parameters-complexity}

The dimensional structure of an FFN determines its parameter count.
Ignoring bias parameters, the input-projection matrix contains
\(d_{\mathrm{ff}}d_{\mathrm{model}}\) entries, and the output-projection
matrix contains \(d_{\mathrm{model}}d_{\mathrm{ff}}\) entries. The total
number of weight parameters is therefore
\begin{align}
    P_{\mathrm{FFN}}
    &=
    d_{\mathrm{ff}}d_{\mathrm{model}}
    +d_{\mathrm{model}}d_{\mathrm{ff}}\\
    &=
    2d_{\mathrm{model}}d_{\mathrm{ff}}.
    \label{eq:ffn-parameter-count}
\end{align}
When both bias vectors are included, the total becomes
\begin{equation}
    P_{\mathrm{FFN,bias}}
    =
    2d_{\mathrm{model}}d_{\mathrm{ff}}
    +d_{\mathrm{ff}}+d_{\mathrm{model}}.
    \label{eq:ffn-parameter-count-bias}
\end{equation}

It is useful to express the intermediate dimension relative to the model
dimension. Define the expansion ratio
\begin{equation}
    r
    =
    \frac{d_{\mathrm{ff}}}{d_{\mathrm{model}}},
    \qquad
    d_{\mathrm{ff}}=r\,d_{\mathrm{model}}.
    \label{eq:ffn-expansion-ratio}
\end{equation}
Substitution into \cref{eq:ffn-parameter-count} gives
\begin{equation}
    P_{\mathrm{FFN}}
    =
    2r\,d_{\mathrm{model}}^2.
    \label{eq:ffn-parameter-expansion}
\end{equation}
For a fixed expansion ratio, the number of FFN weight parameters therefore
grows quadratically with the model hidden dimension. The intermediate
dimension is an architectural choice rather than a universal constant; its
value depends on the model configuration.

The dominant arithmetic cost is incurred by the two matrix
multiplications. Let \(M\) denote the total number of token vectors
processed by the FFN in a batch. The first projection requires
\(M d_{\mathrm{model}}d_{\mathrm{ff}}\) multiply--accumulate operations
(MACs), and the second projection requires the same number. Thus,
\begin{equation}
    C_{\mathrm{FFN}}
    =
    2M d_{\mathrm{model}}d_{\mathrm{ff}}
    \quad\text{MACs}.
    \label{eq:ffn-macs}
\end{equation}
When a multiplication and an addition are counted as separate
floating-point operations, the corresponding leading-order arithmetic
count is
\begin{equation}
    F_{\mathrm{FFN}}
    \approx
    4M d_{\mathrm{model}}d_{\mathrm{ff}}
    \quad\text{FLOPs}.
    \label{eq:ffn-flops}
\end{equation}
This estimate counts the two dominant projections. Bias additions and
element-wise activation operations add work of order
\(O(Md_{\mathrm{ff}})\), which is lower order than the matrix
multiplications when the model dimensions are large. The estimate also
uses the conventional dense-matrix arithmetic count; actual latency
depends on hardware, precision, memory movement, and kernel efficiency.

The distinction between parameter count and arithmetic cost is important.
The matrices are shared across all token positions, so changing the
sequence length does not change the number of FFN parameters. It does,
however, change the number of times those parameters are used and
therefore changes the computation approximately linearly in \(M\).
Increasing \(d_{\mathrm{model}}\) or \(d_{\mathrm{ff}}\), in contrast,
changes both parameter count and per-token computation.

\subsection{Gated Feed-Forward Networks and Their Internal Units}
\label{subsec:ffn-gated}

The two-projection FFN is a useful basic model, but modern language-model
FFNs can introduce an additional projection that controls the
intermediate features. Lo et al. examine the three-projection expert
structure used in LLaMA-style MoE models and write its computation as
\parencite{lo2025closer}:
\begin{equation}
    E(\mathbf{x})
    =
    \mathbf{W}_{\mathrm{down}}
    \left[
        \left(\mathbf{W}_{\mathrm{up}}\mathbf{x}\right)
        \odot
        \sigma\left(\mathbf{W}_{\mathrm{act}}\mathbf{x}\right)
    \right].
    \label{eq:ffn-gated}
\end{equation}
Here, \(\mathbf{x}\in\mathbb{R}^{d_{\mathrm{hid}}}\),
\(\mathbf{W}_{\mathrm{up}},\mathbf{W}_{\mathrm{act}}
\in\mathbb{R}^{d_{\mathrm{mid}}\times d_{\mathrm{hid}}}\), and
\(\mathbf{W}_{\mathrm{down}}
\in\mathbb{R}^{d_{\mathrm{hid}}\times d_{\mathrm{mid}}}\).
The symbol \(\odot\) denotes element-wise multiplication, and
\(\sigma\) is an element-wise activation function. Both the input and
output of the expert belong to \(\mathbb{R}^{d_{\mathrm{hid}}}\), while
the two intermediate projections produce vectors in
\(\mathbb{R}^{d_{\mathrm{mid}}}\).

The up projection produces
\begin{equation}
    \mathbf{a}
    =
    \mathbf{W}_{\mathrm{up}}\mathbf{x},
    \label{eq:ffn-gated-up}
\end{equation}
whereas the activation projection produces
\begin{equation}
    \mathbf{g}
    =
    \sigma\left(\mathbf{W}_{\mathrm{act}}\mathbf{x}\right).
    \label{eq:ffn-gated-control}
\end{equation}
Their element-wise product is
\begin{equation}
    \mathbf{z}
    =
    \mathbf{a}\odot\mathbf{g},
    \qquad
    z_i=a_i g_i,
    \quad i=1,\ldots,d_{\mathrm{mid}}.
    \label{eq:ffn-gated-product}
\end{equation}
The down projection then maps this modulated intermediate representation
back to the hidden dimension:
\begin{equation}
    E(\mathbf{x})
    =
    \mathbf{W}_{\mathrm{down}}\mathbf{z}.
    \label{eq:ffn-gated-down}
\end{equation}
The additional projection is not merely a different choice of activation
function. It supplies a second learned transformation of the input, and
its output controls the intermediate features through element-wise
multiplication.

Ignoring biases, the three matrices contain
\begin{align}
    P_{\mathrm{gated}}
    &=
    d_{\mathrm{mid}}d_{\mathrm{hid}}
    +d_{\mathrm{mid}}d_{\mathrm{hid}}
    +d_{\mathrm{hid}}d_{\mathrm{mid}}\\
    &=
    3d_{\mathrm{hid}}d_{\mathrm{mid}}
    \label{eq:ffn-gated-parameters}
\end{align}
parameters. This is 50\% more projection parameters than a two-projection
FFN with the same input, output, and intermediate dimensions. For
\(M\) tokens, the three projections require approximately
\(3M d_{\mathrm{hid}}d_{\mathrm{mid}}\) MACs, or
\(6M d_{\mathrm{hid}}d_{\mathrm{mid}}\) FLOPs when each
multiply--accumulate is counted as two FLOPs, in addition to
element-wise operations.

The dimensions reported by Lo et al. for the Mistral 7B configuration
include \(d_{\mathrm{hid}}=4096\) and \(d_{\mathrm{mid}}=14336\).
Using these dimensions in the three-projection expression gives
\begin{align}
    P_{\mathrm{gated}}
    &=
    3(4096)(14336)\\
    &=
    176\,160\,768.
    \label{eq:ffn-gated-numeric}
\end{align}
This is the weight count for the three projection matrices under the
bias-free formulation. It is not a count of every parameter in the
Transformer block or the complete language model
\parencite{lo2025closer}.

The three-projection formulation also permits the contribution of each
intermediate unit to be examined explicitly. Let
\(\mathbf{W}_{\mathrm{up}}[i,:]\) and
\(\mathbf{W}_{\mathrm{act}}[i,:]\) denote the \(i\)-th rows of the
up and activation matrices, and let
\(\mathbf{W}_{\mathrm{down}}[:,i]\) denote the \(i\)-th column of the
down matrix. The contribution of unit \(i\) is determined by
\begin{equation}
    z_i
    =
    \left(\mathbf{W}_{\mathrm{up}}[i,:]\mathbf{x}\right)
    \sigma\left(\mathbf{W}_{\mathrm{act}}[i,:]\mathbf{x}\right).
    \label{eq:ffn-neuron-activation}
\end{equation}
The complete expert output can consequently be decomposed as
\begin{equation}
    E(\mathbf{x})
    =
    \sum_{i=1}^{d_{\mathrm{mid}}}
    \mathbf{W}_{\mathrm{down}}[:,i]\,
    \left[
        \left(\mathbf{W}_{\mathrm{up}}[i,:]\mathbf{x}\right)
        \sigma\left(\mathbf{W}_{\mathrm{act}}[i,:]\mathbf{x}\right)
    \right].
    \label{eq:ffn-neuron-decomposition}
\end{equation}
Each intermediate unit therefore combines two input-side projections
with an output-side vector. The input projections determine the
input-dependent scalar activation of that unit, and the corresponding
column of the down matrix determines the direction of its contribution
in the output space.

Lo et al. use this decomposition to motivate the interpretation of FFN
neurons as fine-grained computational experts. This is a functional
interpretation of the internal units, not a claim that each neuron is an
independent expert network with its own router. The activation projection
controls intermediate neurons within one FFN, whereas an MoE router
selects among complete expert modules. These are distinct levels of
computation \parencite{lo2025closer}.

\begin{figure}[htbp]
    \centering
    % Original vector-style schematic for the FFN section.
% This file is input inside a figure environment in
% chapters/02_architecture/03_feed_forward.tex.
% It uses the standard LaTeX picture environment and requires no
% additional drawing package.

\begingroup
\setlength{\unitlength}{0.8pt}
\begin{picture}(420,218)

    % (a) Standard two-projection position-wise FFN
    \put(0,198){\makebox(420,12)[l]{\small\textbf{(a) Standard position-wise FFN}}}
    \put(0,168){\framebox(44,22){\(\mathbf{x}_t\)}}
    \put(76,168){\framebox(68,22){\(\mathbf{W}_1\)}}
    \put(176,168){\framebox(42,22){\(\phi\)}}
    \put(248,168){\framebox(74,22){\(\mathbf{W}_2\)}}
    \put(356,168){\framebox(44,22){\(\mathbf{y}_t\)}}
    \put(44,179){\vector(1,0){32}}
    \put(144,179){\vector(1,0){32}}
    \put(218,179){\vector(1,0){30}}
    \put(322,179){\vector(1,0){34}}

    % (b) Three-projection gated FFN
    \put(0,130){\makebox(420,12)[l]{\small\textbf{(b) Gated three-projection FFN}}}
    \put(0,76){\framebox(44,22){\(\mathbf{x}\)}}
    \put(82,105){\framebox(68,22){\(\mathbf{W}_{\mathrm{up}}\)}}
    \put(82,47){\framebox(68,22){\(\mathbf{W}_{\mathrm{act}}\)}}
    \put(174,47){\framebox(44,22){\(\sigma\)}}
    \put(238,76){\framebox(32,22){\(\odot\)}}
    \put(300,76){\framebox(70,22){\(\mathbf{W}_{\mathrm{down}}\)}}
    \put(390,76){\framebox(30,22){\(E\)}}

    \put(44,87){\vector(4,3){38}}
    \put(44,87){\vector(4,-3){38}}
    \put(150,116){\vector(3,-1){87}}
    \put(150,58){\vector(1,0){24}}
    \put(218,58){\vector(2,1){20}}
    \put(270,87){\vector(1,0){30}}
    \put(370,87){\vector(1,0){20}}

    % (c) Sparse mixture of FFN experts
    \put(0,28){\makebox(420,12)[l]{\small\textbf{(c) Sparse mixture-of-experts FFN}}}
    \put(0,0){\framebox(40,20){\(\mathbf{x}\)}}
    \put(60,0){\framebox(70,20){Router}}
    \put(150,0){\framebox(55,20){Top-\(K\)}}
    \put(225,0){\framebox(80,20){Selected FFNs}}
    \put(325,0){\framebox(42,20){Sum}}
    \put(387,0){\framebox(33,20){\(\mathbf{y}\)}}

    \put(40,10){\vector(1,0){20}}
    \put(130,10){\vector(1,0){20}}
    \put(205,10){\vector(1,0){20}}
    \put(305,10){\vector(1,0){20}}
    \put(367,10){\vector(1,0){20}}

\end{picture}
\endgroup

    \caption{Schematic comparison of a standard position-wise FFN, a
    three-projection gated FFN, and a sparsely routed MoE layer. In the
    standard FFN, the input is projected, transformed element-wise, and
    projected back to the model dimension. In the gated FFN, parallel up
    and activation projections are combined element-wise before the
    down projection. In the sparse MoE layer, a router selects a subset
    of FFN experts whose weighted outputs are aggregated.}
    \label{fig:ffn-architectures}
\end{figure}
% Alt text for accessibility register: Three horizontal computational
% diagrams compare a two-projection FFN, a gated three-projection FFN with
% parallel up and activation branches, and a sparse MoE layer with a router,
% selected experts, and weighted output aggregation.

\subsection{Feed-Forward Networks Within Transformer Blocks}
\label{subsec:ffn-transformer-block}

An FFN is a component of a larger Transformer block rather than an
independent sequence-processing stage. In the formulation used by Dai
et al., the self-attention sublayer first produces a representation with
a residual connection, after which the FFN transforms each token
representation and adds another residual connection
\parencite{dai2024deepseekmoe}. Let
\(\mathbf{h}_{1:T}^{\,l-1}\) denote the token representations entering
block \(l\), and let \(\mathbf{u}_{1:T}^{\,l}\) denote the
representations after its self-attention sublayer. The simplified
computation is
\begin{equation}
    \mathbf{u}_{1:T}^{\,l}
    =
    \operatorname{SelfAtt}
    \left(\mathbf{h}_{1:T}^{\,l-1}\right)
    +
    \mathbf{h}_{1:T}^{\,l-1}.
    \label{eq:ffn-attention-residual}
\end{equation}
The FFN then operates on each token separately:
\begin{equation}
    \mathbf{h}_{t}^{\,l}
    =
    \operatorname{FFN}
    \left(\mathbf{u}_{t}^{\,l}\right)
    +
    \mathbf{u}_{t}^{\,l}.
    \label{eq:ffn-block-residual}
\end{equation}
The residual addition requires the FFN output to have the same dimension
as its input. This is why the second FFN projection returns the
intermediate representation to \(d_{\mathrm{model}}\) dimensions.

A residual connection adds the input of a sublayer to the transformation
computed by that sublayer. It therefore preserves a direct path for the
incoming representation while allowing the learned transformation to
contribute an update. The equations above follow the simplified
formulation in Dai et al.; that paper omits layer-normalization operations
from these equations for brevity. The equations should therefore be read
as a description of the FFN and residual structure, not as a complete
specification of every normalization convention used by all Transformer
implementations \parencite{dai2024deepseekmoe}.

The two sublayers have complementary roles. Self-attention explicitly
combines information across token positions, whereas the position-wise
FFN applies the same learned nonlinear feature transformation to each
token. The FFN does not itself create cross-token interactions; it
transforms the representation that is available at its input.

\subsection{From Dense FFNs to Mixture-of-Experts Layers}
\label{subsec:ffn-to-moe}

The parameter count of an FFN can be substantial because both projection
matrices scale with the model and intermediate dimensions. Cai et al.
report that, in the 540-billion-parameter PaLM configuration discussed
in their survey, approximately 90\% of the parameters are located in
FFN layers. They also summarize an empirical analysis reporting lower
active expert engagement in FFN layers than in self-attention layers
for the particular model and task examined. These observations motivate
the investigation of conditional computation in FFN modules, but the
reported proportions are specific to the cited experiments and should
not be treated as universal constants for Transformer models
\parencite{cai2025survey}.

A mixture-of-experts (MoE) layer replaces one FFN with a collection of
\(N\) expert networks,
\[
    \mathcal{E}
    =
    \{E_1,E_2,\ldots,E_N\},
\]
where each expert is itself an FFN. A gating network, also called a
router, assigns scores to the experts for each input token. In a
linear-router formulation, the router logits for input
\(\mathbf{x}\in\mathbb{R}^{d_{\mathrm{model}}}\) are
\begin{equation}
    \mathbf{q}(\mathbf{x})
    =
    \mathbf{W}_{r}\mathbf{x},
    \qquad
    \mathbf{W}_{r}
    \in\mathbb{R}^{N\times d_{\mathrm{model}}}.
    \label{eq:moe-router-logits}
\end{equation}
The logits are transformed into routing probabilities by the softmax:
\begin{equation}
    p_i(\mathbf{x})
    =
    \frac{\exp(q_i(\mathbf{x}))}
    {\displaystyle\sum_{j=1}^{N}\exp(q_j(\mathbf{x}))},
    \qquad
    \sum_{i=1}^{N}p_i(\mathbf{x})=1.
    \label{eq:moe-router-softmax}
\end{equation}
The probabilities are nonnegative and sum to one, so they define a
distribution over the available experts.

In a dense MoE layer, all experts are evaluated and their outputs are
combined using their routing probabilities:
\begin{equation}
    \operatorname{MoE}_{\mathrm{dense}}(\mathbf{x})
    =
    \sum_{i=1}^{N}p_i(\mathbf{x})E_i(\mathbf{x}).
    \label{eq:moe-dense}
\end{equation}
In a sparse MoE layer, only the \(K\) experts with the largest routing
scores are evaluated. Let \(\mathcal{T}_K(\mathbf{x})\) denote the set
of their indices. The sparse gate is
\begin{equation}
    g_i(\mathbf{x})
    =
    \begin{cases}
        p_i(\mathbf{x}),
        & i\in\mathcal{T}_K(\mathbf{x}),\\
        0,
        & i\notin\mathcal{T}_K(\mathbf{x}).
    \end{cases}
    \label{eq:moe-sparse-gate}
\end{equation}
The sparse MoE output is
\begin{equation}
    \operatorname{MoE}_{\mathrm{sparse}}(\mathbf{x})
    =
    \sum_{i=1}^{N}g_i(\mathbf{x})E_i(\mathbf{x})
    =
    \sum_{i\in\mathcal{T}_K(\mathbf{x})}
    p_i(\mathbf{x})E_i(\mathbf{x}).
    \label{eq:moe-sparse-output}
\end{equation}
The distinction is computational: in the sparse formulation, experts
whose gate values are zero need not perform their FFN computation for
that token. The linear-router formulation above follows the routing
description in Fedus et al.; Dai et al. also describe token-to-expert
affinity scores based on the similarity between a token representation
and an expert centroid \parencite{fedus2022switch,dai2024deepseekmoe}.

If each expert contains \(P_E\) parameters and all experts have the same
size, the total expert parameter count is
\begin{equation}
    P_{\mathrm{experts,total}}
    =
    NP_E.
    \label{eq:moe-total-expert-parameters}
\end{equation}
When \(K\) experts are active for a token, the number of active expert
parameters is
\begin{equation}
    P_{\mathrm{experts,active}}
    =
    KP_E.
    \label{eq:moe-active-expert-parameters}
\end{equation}
The ratio is therefore
\begin{equation}
    \frac{P_{\mathrm{experts,active}}}
    {P_{\mathrm{experts,total}}}
    =
    \frac{K}{N}.
    \label{eq:moe-active-ratio}
\end{equation}
These expressions concern expert parameters and exclude router
parameters and any other shared components. They show how a sparse
MoE layer can expose a large total pool of FFN parameters while
activating only a subset for an individual token. The total model
capacity and the expert computation per token are consequently
different quantities.

\subsection{Switch Routing, Expert Capacity, and Load Balancing}
\label{subsec:ffn-switch-routing}

Fedus et al. introduce the Switch Transformer as a simplified sparse
MoE design in which each token is routed to one expert rather than to
multiple experts. The selected expert is
\begin{equation}
    i^{*}
    =
    \operatorname*{arg\,max}_{1\leq i\leq N}
    p_i(\mathbf{x}),
    \label{eq:switch-top-one}
\end{equation}
and the output of the Switch FFN is
\begin{equation}
    \operatorname{SwitchFFN}(\mathbf{x})
    =
    p_{i^{*}}(\mathbf{x})E_{i^{*}}(\mathbf{x}).
    \label{eq:switch-output}
\end{equation}
The expert output is scaled by the selected router probability. Top-1
routing reduces the number of expert computations per token and
simplifies token dispatch relative to a routing rule that activates
multiple experts. Fedus et al. report that this simplification can
preserve model quality in their experiments while reducing routing
and communication requirements \parencite{fedus2022switch}.

Sparse routing creates a finite-capacity allocation problem. Let \(M\)
be the number of tokens presented to one MoE layer in a batch. Under
uniform allocation across \(N\) experts, the expected number of tokens
per expert is \(M/N\). Switch Transformer defines expert capacity as
\begin{equation}
    C
    =
    \left(\frac{M}{N}\right)c_{\mathrm{factor}},
    \label{eq:switch-expert-capacity}
\end{equation}
where \(c_{\mathrm{factor}}\) is the capacity factor. A factor greater
than one provides buffer for nonuniform routing. In an implementation,
the resulting capacity must be represented by an admissible integer
tensor dimension.

For example, consider \(M=1024\) tokens and \(N=16\) experts. The
uniform allocation is \(1024/16=64\) tokens per expert. With
\(c_{\mathrm{factor}}=1.5\), the capacity is
\[
    C=64\times1.5=96.
\]
An expert can therefore accept up to 96 tokens in this example.
If more tokens are routed to it, the excess tokens are dropped from
that expert computation in the Switch implementation and their
representations are passed to the next layer through the residual
connection. A larger capacity factor reduces the risk of overflow,
but it also allocates more buffer and can increase memory and
computation spent on unused capacity. Capacity is therefore a
trade-off between token overflow and resource utilization
\parencite{fedus2022switch}.

The router may not distribute tokens uniformly. To encourage a more
balanced allocation, Switch Transformer adds an auxiliary loss to the
main training objective. For a batch \(\mathcal{B}\) containing \(M\)
tokens, define the fraction of tokens dispatched to expert \(i\) as
\begin{equation}
    f_i
    =
    \frac{1}{M}
    \sum_{\mathbf{x}\in\mathcal{B}}
    \mathbb{1}
    \left[
        \operatorname*{arg\,max}_{j}p_j(\mathbf{x})=i
    \right],
    \label{eq:switch-dispatch-fraction}
\end{equation}
where \(\mathbb{1}[\cdot]\) is the indicator function. Define the
average router probability assigned to expert \(i\) as
\begin{equation}
    P_i
    =
    \frac{1}{M}
    \sum_{\mathbf{x}\in\mathcal{B}}p_i(\mathbf{x}).
    \label{eq:switch-router-fraction}
\end{equation}
The auxiliary load-balancing loss is
\begin{equation}
    \mathcal{L}_{\mathrm{aux}}
    =
    \alpha N\sum_{i=1}^{N}f_iP_i,
    \label{eq:switch-load-balancing}
\end{equation}
where \(\alpha\) controls the contribution of this term to the training
objective. The vector \(f\) records the realized dispatch frequencies,
whereas \(P\) records the average probability mass assigned by the
router. These quantities are not interchangeable: one describes
discrete assignments and the other describes the router's probability
outputs.

Uniform allocation corresponds to \(f_i=P_i=1/N\) for every expert.
Under this condition,
\[
    \mathcal{L}_{\mathrm{aux}}
    =
    \alpha N
    \sum_{i=1}^{N}\frac{1}{N^2}
    =
    \alpha.
\]
The objective is minimized at the uniform allocation. Although the
dispatch indicator in \(f_i\) is discrete, the probability term \(P_i\)
depends on the softmax router outputs and supplies a differentiable
training signal to the router. The auxiliary loss is added to the
model's main loss; it does not replace the language-modeling
objective \parencite{fedus2022switch}.

\subsection{Fine-Grained Segmentation and Shared Experts}
\label{subsec:ffn-deepseekmoe}

A conventional sparse MoE layer may contain a limited number of
relatively large FFN experts. Dai et al. identify two potential
limitations of this organization. Knowledge hybridity occurs when
one expert must process tokens associated with diverse types of
knowledge. Knowledge redundancy occurs when different experts
independently acquire overlapping information that is useful across
multiple contexts. DeepSeekMoE addresses these problems through
fine-grained expert segmentation and shared expert isolation
\parencite{dai2024deepseekmoe}.

In fine-grained expert segmentation, each original expert is divided
into \(m\) smaller experts by reducing its intermediate dimension by
a factor of \(m\). If the original architecture has \(N\) experts,
each with intermediate dimension \(d_{\mathrm{ff}}\), the fine-grained
architecture has
\begin{equation}
    N_{\mathrm{fine}}=mN,
    \qquad
    d_{\mathrm{ff,fine}}=\frac{d_{\mathrm{ff}}}{m}.
    \label{eq:deepseek-fine-grained-dimensions}
\end{equation}
To preserve the dominant expert computation per token, the number of
activated experts is increased from \(K\) to \(mK\):
\begin{equation}
    K_{\mathrm{fine}}=mK.
    \label{eq:deepseek-fine-grained-topk}
\end{equation}
These relations assume that the intermediate dimension can be divided
as specified and that the expert projection costs scale linearly with
that dimension.

If an original expert contains \(P_E\) projection parameters, each
fine-grained expert contains approximately \(P_E/m\) such parameters.
The total expert parameter count is therefore preserved:
\begin{align}
    P_{\mathrm{fine,total}}
    &=
    mN\left(\frac{P_E}{m}\right)\\
    &=
    NP_E.
    \label{eq:deepseek-total-parameters}
\end{align}
Likewise, if the dominant computation of one original expert is
\(C_E\), then the active expert computation is
\begin{align}
    C_{\mathrm{fine,active}}
    &=
    mK\left(\frac{C_E}{m}\right)\\
    &=
    KC_E.
    \label{eq:deepseek-active-computation}
\end{align}
Thus, the segmentation increases the number of available experts and
the number selected for each token while preserving the leading-order
expert parameter count and active expert computation.

The principal change is the number of possible expert combinations.
If \(K\) experts are selected from \(N\) experts, the number of
unordered combinations is
\begin{equation}
    \binom{N}{K}
    =
    \frac{N!}{K!(N-K)!}.
    \label{eq:moe-combination-count}
\end{equation}
After segmentation, the corresponding count is
\begin{equation}
    \binom{mN}{mK}.
    \label{eq:deepseek-combination-count}
\end{equation}
Dai et al. give the example \(N=16\) and \(K=2\), for which the
conventional architecture permits
\[
    \binom{16}{2}=120
\]
possible combinations. With \(m=4\), the fine-grained architecture
contains 64 experts and activates eight, giving
\[
    \binom{64}{8}
    =
    4\,426\,165\,368
\]
possible combinations. This is a count of possible subsets, not a
claim that all combinations are equally likely or are observed during
training. The larger combination space gives the router more ways to
compose specialized expert subsets \parencite{dai2024deepseekmoe}.

Fine-grained segmentation does not, by itself, eliminate redundancy
between experts. DeepSeekMoE therefore introduces shared expert
isolation. A designated set of \(K_s\) shared experts is activated for
every token, while the remaining experts are routed conditionally.
The number of routed experts is \(mN-K_s\), and the number of routed
experts selected for each token is \(mK-K_s\). A simplified
representation of the layer output is
\begin{align}
    \mathbf{h}_{t}^{\,l}
    ={}&
    \sum_{j=1}^{K_s}
    E^{\mathrm{shared}}_j
    \left(\mathbf{u}_{t}^{\,l}\right)
    \nonumber\\
    &+
    \sum_{i=1}^{mN-K_s}
    g_{i,t}
    E^{\mathrm{routed}}_i
    \left(\mathbf{u}_{t}^{\,l}\right)
    +
    \mathbf{u}_{t}^{\,l}.
    \label{eq:deepseek-shared-routed}
\end{align}
The shared experts provide a common computational pathway for
information that is useful across contexts. The routed experts can
then specialize more strongly in input-dependent information without
each having to reproduce all common knowledge.

For the routing formulation described by Dai et al., the affinity
between token \(t\) and expert \(i\) in layer \(l\) is computed from
the token representation and the expert centroid
\(\mathbf{e}_{i}^{\,l}\):
\begin{equation}
    s_{i,t}
    =
    \operatorname{Softmax}_{i}
    \left(
        (\mathbf{u}_{t}^{\,l})^{\top}\mathbf{e}_{i}^{\,l}
    \right).
    \label{eq:deepseek-centroid-affinity}
\end{equation}
The sparse gate retains the \(mK\) largest affinities in the
fine-grained formulation, with shared experts handled separately.
The centroid-affinity expression is a particular routing
parameterization and should not be conflated with the linear router
used in the generic Switch formulation.

The DeepSeekMoE paper distinguishes balancing at the expert level
from balancing at the device level. Let
\[
    N'=mN-K_s,
    \qquad
    K'=mK-K_s
\]
denote the number of routed experts and the number of routed experts
selected for each token, respectively. For a batch of \(T\) tokens,
the normalized dispatch fraction for routed expert \(i\) is
\begin{equation}
    f_i
    =
    \frac{N'}{K'T}
    \sum_{t=1}^{T}
    \mathbb{1}
    \left(\text{token }t\text{ selects routed expert }i\right),
    \label{eq:deepseek-expert-dispatch}
\end{equation}
and the average affinity assigned to that expert is
\begin{equation}
    P_i
    =
    \frac{1}{T}\sum_{t=1}^{T}s_{i,t}.
    \label{eq:deepseek-expert-affinity}
\end{equation}
The expert-level balancing objective is
\begin{equation}
    \mathcal{L}_{\mathrm{ExpBal}}
    =
    \frac{\alpha_1}{N'}
    \sum_{i=1}^{N'}f_iP_i,
    \label{eq:deepseek-expert-balance}
\end{equation}
where \(\alpha_1\) controls the strength of expert-level balancing.
The normalization in \(f_i\) makes its value equal to one under
uniform routed-expert allocation.

Balancing individual experts does not guarantee balanced accelerator
workloads. Suppose the routed experts are partitioned into \(D\) groups,
with group \(\mathcal{G}_d\) deployed on device \(d\). The normalized
dispatch fraction and aggregate affinity for device group \(d\) are
\begin{align}
    f'_d
    &=
    \frac{1}{|\mathcal{G}_d|}
    \sum_{j\in\mathcal{G}_d}f_j,
    \label{eq:deepseek-device-dispatch}\\
    P'_d
    &=
    \sum_{j\in\mathcal{G}_d}P_j.
    \label{eq:deepseek-device-affinity}
\end{align}
The corresponding device-level balancing objective is
\begin{equation}
    \mathcal{L}_{\mathrm{DevBal}}
    =
    \alpha_2
    \sum_{d=1}^{D}f'_dP'_d,
    \label{eq:deepseek-device-balance}
\end{equation}
where \(\alpha_2\) is the device-level balance factor. Dai et al.
describe using a relatively small expert-level balance factor to
reduce routing collapse while assigning a larger weight to
device-level balancing, because overly strict expert-level balancing
can compromise model performance. These two objectives address
different bottlenecks: one concerns token allocation among experts,
and the other concerns the aggregate computational load assigned
to devices \parencite{dai2024deepseekmoe}.

Dai et al. report that DeepSeekMoE 2B achieves performance comparable
to GShard 2.9B despite the latter having approximately 1.5 times the
expert parameters and computation. They also report that DeepSeekMoE
16B achieves performance comparable to LLaMA2 7B using about 40\%
of the computation in their reported comparison. These are empirical
results for the configurations and evaluations in that paper, rather
than general guarantees for every fine-grained MoE design
\parencite{dai2024deepseekmoe}.

\subsection{Expert Structure, Specialization, and Practical Constraints}
\label{subsec:ffn-expert-analysis}

The presence of multiple experts does not by itself establish that
they have learned distinct functions. Experts may acquire overlapping
parameters or exhibit similar output behavior. Lo et al. investigate
this issue by analyzing the projection matrices
\(\mathbf{W}_{\mathrm{up}}\),
\(\mathbf{W}_{\mathrm{act}}\), and
\(\mathbf{W}_{\mathrm{down}}\) in several MoE models
\parencite{lo2025closer}.

One way to compare two parameter matrices is to flatten them into
vectors and calculate cosine similarity. For nonzero vectors
\(\mathbf{a},\mathbf{b}\in\mathbb{R}^{n}\), cosine similarity is
\begin{equation}
    \operatorname{CosSim}(\mathbf{a},\mathbf{b})
    =
    \frac{\mathbf{a}^{\top}\mathbf{b}}
    {\|\mathbf{a}\|_2\|\mathbf{b}\|_2},
    \label{eq:ffn-cosine-similarity}
\end{equation}
where
\[
    \|\mathbf{a}\|_2
    =
    \sqrt{\sum_{j=1}^{n}a_j^2}
\]
is the Euclidean norm. Cosine similarity measures angular alignment
and is invariant to positive rescaling of either vector. It does not,
however, measure equality of magnitudes or guarantee equality of
the functions implemented by the corresponding networks.

Flattened-matrix comparison is also sensitive to neuron ordering.
Two experts may contain similar units in different orders, while a
direct element-by-element comparison pairs units that do not
correspond. Lo et al. therefore investigate neuron-level comparisons,
including a reordering procedure based on the Jonker--Volgenant
algorithm for a linear assignment problem. Their analysis treats a
gated FFN neuron as the combination of the corresponding rows of
\(\mathbf{W}_{\mathrm{up}}\) and \(\mathbf{W}_{\mathrm{act}}\), together
with the corresponding column of \(\mathbf{W}_{\mathrm{down}}\).
This connects the neuron decomposition in
\cref{eq:ffn-neuron-decomposition} to a more detailed analysis of
expert similarity \parencite{lo2025closer}.

In their investigated models, Lo et al. report that matrix-level
similarities differ across architectures and layers. For Mixtral,
they report cosine similarities primarily in the range 0.2--0.4
between experts, while similarities between Mixtral experts and the
dense Mistral FFN are approximately 0.6 in the examined comparisons.
They also report that similarities between the routed experts of
DeepSeekMoE and the experts of Grok-1 are close to zero in their
analysis. These observations are specific to the models, layers,
matrices, and comparison procedures studied. They should not be
interpreted as a universal numerical measure of expert
specialization.

The choice of expert count, expert dimension, number of active
experts, and frequency of MoE layers also affects the architecture.
Cai et al. identify these as important design variables and describe
configurations in which MoE layers replace every second FFN layer,
every fourth FFN layer, or the FFN in every layer. A larger number
of experts can increase total parameter capacity, while a larger
expert dimension increases the capacity and cost of each expert.
Increasing the number of active experts changes the computation
performed per token. These choices must therefore be considered
together rather than optimized independently
\parencite{cai2025survey}.

The theoretical arithmetic savings of sparse FFNs do not automatically
translate into proportional reductions in wall-clock latency. Routing
requires score calculation and expert selection. The implementation
must also dispatch tokens to the selected experts and combine their
outputs. When experts are distributed across accelerators, this can
introduce communication and synchronization costs. Zhang et al.
review practical issues including irregular memory access,
cross-device communication, fragmented workloads, and hardware
underutilization. These system effects can reduce the practical
benefit of sparse computation, particularly under latency or memory
constraints \parencite{zhang2025moe}.

The architecture must consequently balance representational capacity,
expert specialization, routing quality, load distribution, memory
requirements, and the communication pattern imposed by the hardware.
An MoE layer is not efficient merely because it activates a small
fraction of its parameters; its routing and execution mechanisms must
also use the available computational resources effectively.

\subsection{Limitations and Connection to Activation Functions}
\label{subsec:ffn-limitations-transition}

The development of FFN architectures reveals a progression from a
shared nonlinear transformation applied to every token, through
gated intermediate representations, to conditional computation
over multiple FFN experts. The conventional FFN provides a
position-wise mapping whose parameter count and arithmetic cost are
determined by its input, intermediate, and output dimensions. Gated
FFNs introduce an additional learned projection that modulates
intermediate features. MoE architectures extend the same basic
feed-forward computation by selecting a subset of expert networks
for each token.

These architectures introduce different trade-offs. Increasing the
intermediate dimension increases the capacity and cost of a dense
FFN. Adding a gate projection increases the parameterization and
computation of a gated FFN. Increasing the number of MoE experts
increases total expert capacity, but sparse execution introduces
routing, capacity, load-balancing, and communication constraints.
Fine-grained segmentation and shared expert isolation provide
strategies for organizing expert parameters, but they do not remove
the need to evaluate the quality of specialization or the practical
cost of routing.

The mathematical role of the activation function remains central
throughout this progression. In a conventional FFN, it supplies the
nonlinearity between the two projections. In a gated FFN, an
activation function transforms the control projection before it is
combined with the up-projection. In an MoE architecture, the router
also uses a probability transformation to produce expert-selection
weights. These operations may all involve nonlinear functions, but
they act on different representations and serve different
architectural purposes. The next section examines activation
functions as mathematical operators, establishing the properties
needed to understand their effects on feature transformations,
gated computation, and gradient-based learning.
